\begin{table}[htbp]
\centering
\caption{Heterogeneous Effects on Nightlight Growth}
\label{tab:heterogeneity}
\begin{tabular}{l ccccc}
\toprule
Subgroup & Estimate & Robust SE & $p$-value & Bandwidth & Eff.\ $N$ \\
\midrule
\multicolumn{6}{l}{\textit{Panel A: Crime Severity}} \\
\addlinespace[3pt]
Major Crimes                   & --- & --- & --- & --- & 17 \\
Minor Crimes Only              & 0.136$^{*}$ & 0.076 & 0.073 & 10.82 & 1,230 \\
\addlinespace[6pt]
\multicolumn{6}{l}{\textit{Panel B: State Development}} \\
\addlinespace[3pt]
BIMARU States                  & 0.294$^{**}$ & 0.143 & 0.040 & 8.57 & 443 \\
Non-BIMARU States              & -0.019 & 0.056 & 0.728 & 14.28 & 883 \\
\addlinespace[6pt]
\multicolumn{6}{l}{\textit{Panel C: Seat Reservation}} \\
\addlinespace[3pt]
General Seats                  & 0.086 & 0.075 & 0.255 & 10.70 & 1,004 \\
Reserved Seats (SC/ST)         & 0.362 & 0.281 & 0.198 & 10.95 & 232 \\
\bottomrule
\end{tabular}
\begin{minipage}{0.92\textwidth}
\vspace{4pt}
\footnotesize
\textit{Notes:} Each row reports a separate RDD regression of nightlight growth on the
criminal-candidate vote margin for the indicated subgroup.
Panel A splits by crime severity (major crimes include murder, kidnapping,
extortion per ADR classification).
Panel B compares BIMARU states (Bihar, MP, Rajasthan, UP, and successor states)
with non-BIMARU states.
Panel C compares General category seats with those reserved for SC/ST.
All specifications use local linear estimation with MSE-optimal bandwidth.
$^{***}$ $p<0.01$, $^{**}$ $p<0.05$, $^{*}$ $p<0.1$.
\end{minipage}
\end{table}
